% Generated from content/cv.md. Edit the Markdown source, then run make cv-tex.
\documentclass[10pt,letterpaper]{article}
\usepackage[letterpaper,margin=0.72in]{geometry}
\usepackage{fontspec}
\setmainfont{Latin Modern Roman}
\usepackage{microtype}
\usepackage[hidelinks]{hyperref}
\usepackage{enumitem}
\usepackage{titlesec}
\usepackage{xcolor}

\pagestyle{empty}
\setlength{\parindent}{0pt}
\setlength{\parskip}{5pt}
\setlist[itemize]{leftmargin=1.2em,itemsep=2pt,topsep=2pt,parsep=0pt}
\titleformat{\section}{\large\bfseries\uppercase}{}{0pt}{}[\titlerule]
\titleformat{\subsection}{\normalsize\bfseries}{}{0pt}{}
\titlespacing*{\section}{0pt}{15pt}{5pt}
\titlespacing*{\subsection}{0pt}{9pt}{1pt}
\newcommand{\cvtitle}[1]{{\Huge\bfseries #1}\vspace{3pt}}

\begin{document}
\cvtitle{Robert A. Cochran III}

Technical engineering leader | security systems, Kubernetes, software verification, AI-agent infrastructure

San Francisco, CA / \href{mailto:kudzoo@gmail.com}{kudzoo@gmail.com} / \href{https://www.robbycochran.com/}{robbycochran.com} / \href{https://github.com/robbycochran}{GitHub} / \href{https://scholar.google.com/citations?user=qW0wkZIAAAAJ\&hl=en}{Google Scholar} / \href{https://www.linkedin.com/in/robertacochran}{LinkedIn}

\section*{Profile}

Technical engineering leader and security systems researcher with a PhD in computer science and a publication record across NDSS, NSDI, POPL, and TISSEC. Work spans symbolic verification of remote client behavior, program synthesis, Kubernetes runtime security, eBPF-based observability, enterprise product architecture, engineering management of a five-engineer C++/Go/Rust team, customer/design-partner strategy, product strategy, and current ACS agentic foundations: secure harness primitives, skills strategy, support triage automation, internal hackathon leadership, and CI-tested agent workflows.

\section*{Interests}

Runtime security and behavioral detection, Kubernetes / OpenShift / RHACS, eBPF observability, Software verification and formal methods, Symbolic execution, Program synthesis, AI-agent safety and secure harnesses, Policy-as-code and detection layering, Developer/security workflows

\section*{Experience and research record}

\subsection*{IBM / Red Hat - Red Hat Advanced Cluster Security (RHACS / ACS) — Staff Engineer / Engineering Manager}

\emph{San Francisco, CA / Remote | 2021 - Present}

\begin{itemize}

  \item Lead runtime-security product and architecture work across RHACS process activity, network security, file activity, behavioral baselines, runtime anomaly detection, Kubernetes-native policy workflows, and analyst/developer experience.

  \item Lead a team of five C++/Go/Rust engineers and balance people management, architecture, product strategy, and hands-on technical direction.

  \item Drive eBPF-based runtime visibility design for OpenShift/Kubernetes environments, including collection semantics, policy exceptions, performance constraints, platform support, and low false-positive detection goals.

  \item Lead ACS agentic foundations: internal hackathon, CI tooling, support triage bots, skills strategy, secure harness primitives, and developer paths from local sandbox to OpenShift execution.

  \item Build and guide secure AI-agent infrastructure patterns through harness-openshell: declarative agent workflows, provider/credential boundaries, inference routing, policy-aware sandbox execution, and repeatable CI validation.

  \item Work with large enterprise customers and design partners to translate deployment constraints into roadmap decisions around scale, compliance, architecture, performance, IBM Power/Z support, and field-support workflows.

  \item Managed distributed engineering teams while remaining close to architecture and implementation; promoted L4 and L5 engineers through scope design, calibrated feedback, and career planning.

  \item Shaped OpenShift-first RHACS strategy around console integration, RBAC alignment, GitOps/policy-as-code, developer experience, runtime observability, and AI-assisted workflows.

  \item Outputs from this period:

  \item \href{https://github.com/stackrox/harness-openshell}{Proof object: harness-openshell - declarative configuration harness for OpenShell AI agent sandboxes.}

  \item \href{https://patents.justia.com/patent/11811804}{Patent: US 11,811,804, "System and method for detecting process anomalies in a distributed computation system utilizing containers"; Red Hat assignee; named inventor.}

  \item \href{https://speakerdeck.com/redhatlivestreaming/stackrox-community-office-hours-e2-ebpf-101-implementing-security-and-monitoring-kubernetes}{Public talk: eBPF 101: implementing security and monitoring in Kubernetes.}

\end{itemize}

\subsection*{StackRox, Inc. (acquired by Red Hat) — Member of Technical Staff / Senior Software Engineer / Staff Software Engineer}

\emph{Mountain View, CA | May 2017 - 2021}

\begin{itemize}

  \item Joined an early-stage Kubernetes security startup and built runtime detection and enforcement features for containerized workloads and Kubernetes clusters.

  \item Advanced core runtime-security capabilities including baseline-driven detection, process and network visibility, and connections between observed runtime behavior and security policy.

  \item Worked across product, engineering, and customer-driven security cases with large design partners, helping mature the platform through acquisition by Red Hat.

  \item Contributed through the StackRox/RHACS product arc from early startup product to enterprise OpenShift security portfolio; product context included \$0 to \$50M+ revenue and 100+ releases.

  \item Outputs from this period:

  \item \href{https://www.redhat.com/en/blog/faq-red-hat-acquire-stackrox}{Company milestone: StackRox acquisition by Red Hat.}

  \item \href{https://www.redhat.com/en/technologies/cloud-computing/openshift/advanced-cluster-security-kubernetes}{Product: Red Hat Advanced Cluster Security for Kubernetes.}

  \item \href{https://www.youtube.com/watch?v=LFqWl7XB5Gc}{Public talk: Listen to your Engine: Unearthing Security Signals from the Modern Linux Kernel. BSidesSF 2018.}

\end{itemize}

\subsection*{University of North Carolina at Chapel Hill — Research Assistant, Computer Security and Software Verification}

\emph{Chapel Hill, NC | 2008 - 2016}

\begin{itemize}

  \item Developed symbolic client verification: server-side methods for checking whether observed client network traffic could have been produced by sanctioned software.

  \item Built and evaluated symbolic-execution systems for online games, cryptographic clients, and distributed applications, including OpenSSL/Heartbleed-style client-misbehavior detection.

  \item Investigated data-driven symbolic execution using instruction traces, clustering, and novel metrics; refactored KLEE symbolic execution internals for thread-level parallelism and many-core execution.

  \item Research outputs included NDSS 2010 Best Paper, NDSS 2013, NSDI 2017, ACM TISSEC 2011, and dissertation work on remote client behavior.

  \item Outputs from this period:

  \item \href{https://cdr.lib.unc.edu/concern/dissertations/p8418p298}{Dissertation: Symbolic Verification of Remote Client Behavior in Distributed Systems.}

  \item \href{https://www.usenix.org/conference/nsdi17/technical-sessions/presentation/chi}{A System to Verify Network Behavior of Known Cryptographic Clients. NSDI 2017.}

  \item \href{https://www.ndss-symposium.org/ndss-2013-programme/toward-online-verification-client-behavior-distributed-applications/}{Toward Online Verification of Client Behavior in Distributed Applications. NDSS 2013.}

  \item \href{https://www.ndss-symposium.org/wp-content/uploads/2017/09/beth.pdf}{Server-side Verification of Client Behavior in Online Games. NDSS 2010 Best Paper; ACM TISSEC 2011 journal version.}

\end{itemize}

\subsection*{Microsoft Research - Research in Software Engineering (RiSE) — Research Intern with Dr. Ben Livshits}

\emph{Redmond, WA | Summer 2013}

\begin{itemize}

  \item Designed a program-synthesis workflow combining symbolic automata, genetic programming, and online workers to improve regular expressions from examples and human feedback.

  \item Outputs from this period:

  \item \href{https://www.microsoft.com/en-us/research/publication/program-boosting-program-synthesis-via-crowd-sourcing/}{Publication: Program Boosting: Program Synthesis via Crowd-Sourcing. POPL 2015.}

  \item \href{https://patents.google.com/patent/US9753696B2}{Patent: US 9,753,696, "Program boosting including using crowdsourcing for correctness"; Microsoft Technology Licensing assignee; named inventor.}

\end{itemize}

\subsection*{Intel Corporation - Visual Computing Group — Software Engineer}

\emph{Hillsboro, OR | 2006 - 2007}

\begin{itemize}

  \item Built graphics-hardware analysis tooling to interpret, record, and replay Direct3D/OpenGL API calls for performance analysis of prototype visual-computing hardware.

\end{itemize}

\subsection*{Clemson University — Undergraduate Researcher}

\emph{Clemson, SC | 2005 - 2006}

\begin{itemize}

  \item Designed and implemented GPU-based rendering research on second-order illumination and built network-simulation tooling for DOCSIS MAC-layer analysis.

  \item Outputs from this period:

  \item \href{https://www.robbycochran.com/pdf/ACMSE07.pdf}{Publication: Second-order Illumination in Real-Time. ACM Southeast 2007; Best Paper.}

\end{itemize}

\section*{Education}

\begin{itemize}

  \item \textbf{PhD, Computer Science}, University of North Carolina at Chapel Hill, 2016.

  \item Thesis: Symbolic Verification of Remote Client Behavior in Distributed Systems

  \item Advisor: Prof. Michael K. Reiter

  \item \textbf{BS, Computer Science, Magna Cum Laude}, Clemson University, 2006.

  \item Outstanding Senior in Computer Science; Donald A. Norton Computer Science Scholarship; Upsilon Pi Epsilon Honor Society

\end{itemize}

\section*{Publications}

\begin{itemize}

  \item \href{https://www.usenix.org/conference/nsdi17/technical-sessions/presentation/chi}{A. Chi, R. Cochran, M. Nesfield, M. K. Reiter, C. Sturton. "A System to Verify Network Behavior of Known Cryptographic Clients." USENIX NSDI, 2017.} \emph{Origin: UNC.}

  \item \href{https://www.microsoft.com/en-us/research/publication/program-boosting-program-synthesis-via-crowd-sourcing/}{R. Cochran, L. D'Antoni, B. Livshits, M. Veanes. "Program Boosting: Program Synthesis via Crowd-Sourcing." POPL, 2015.} \emph{Origin: Microsoft Research.}

  \item \href{https://www.ndss-symposium.org/ndss-2013-programme/toward-online-verification-client-behavior-distributed-applications/}{R. Cochran, M. K. Reiter. "Toward Online Verification of Client Behavior in Distributed Applications." NDSS, 2013.} \emph{Origin: UNC.}

  \item \href{https://doi.org/10.1145/2043628.2043633}{D. Bethea, R. Cochran, M. K. Reiter. "Server-side Verification of Client Behavior in Online Games." ACM TISSEC, 2011.} \emph{Origin: UNC.}

  \item \href{https://www.ndss-symposium.org/wp-content/uploads/2017/09/beth.pdf}{D. Bethea, R. Cochran, M. K. Reiter. "Server-side Verification of Client Behavior in Online Games." NDSS, 2010. Best Paper.} \emph{Origin: UNC.}

  \item \href{https://www.robbycochran.com/pdf/ACMSE07.pdf}{R. Cochran, J. Steele. "Second-order Illumination in Real-Time." ACM Southeast, 2007. Best Paper.} \emph{Origin: Clemson.}

\end{itemize}

\section*{Patents}

\begin{itemize}

  \item \href{https://patents.justia.com/patent/11811804}{US 11,811,804. "System and method for detecting process anomalies in a distributed computation system utilizing containers." Red Hat, Inc.; filed Dec. 15, 2020; issued Nov. 7, 2023. Named inventor.} \emph{Origin: StackRox / Red Hat.}

  \item \href{https://patents.google.com/patent/US9753696B2}{US 9,753,696. "Program boosting including using crowdsourcing for correctness." Microsoft Technology Licensing LLC; filed Mar. 14, 2014; issued Sept. 5, 2017. Named inventor.} \emph{Origin: Microsoft Research.}

\end{itemize}

\section*{Awards}

\begin{itemize}

  \item Best Paper Award, Network and Distributed System Security Symposium (NDSS), 2010.

  \item Best Paper Award, ACM Southeast Regional Conference, 2007.

  \item Outstanding Senior in Computer Science, Clemson University, 2006.

  \item Donald A. Norton Computer Science Scholarship, Clemson University, 2006.

  \item Upsilon Pi Epsilon Honor Society, Clemson University, 2005.

\end{itemize}

\section*{Service}

\begin{itemize}

  \item External reviewer, ACM Conference on Data and Application Security and Privacy, 2013-2014.

  \item External reviewer, Network and Distributed System Security Symposium, 2013-2014.

\end{itemize}

\section*{Talks and presentations}

\begin{itemize}

  \item \href{https://www.youtube.com/watch?v=LFqWl7XB5Gc}{Listen to your Engine: Unearthing Security Signals from the Modern Linux Kernel. BSidesSF, 2018.} (\href{https://allbsides.com/talk/LFqWl7XB5Gc.html}{AllBSides archive})

  \item \href{https://speakerdeck.com/redhatlivestreaming/stackrox-community-office-hours-e2-ebpf-101-implementing-security-and-monitoring-kubernetes}{eBPF 101: implementing security and monitoring in Kubernetes. StackRox / Red Hat community office hours.}

  \item \href{https://www.ndss-symposium.org/ndss-2013-programme/toward-online-verification-client-behavior-distributed-applications/}{Toward Online Verification of Client Behavior in Distributed Applications. NDSS, San Diego, CA, February 2013.}

  \item \href{https://www.ndss-symposium.org/wp-content/uploads/2017/09/beth.pdf}{Server-side Verification of Client Behavior in Online Games. NDSS, San Diego, CA, February 2010.}

  \item \href{https://www.robbycochran.com/pdf/ACMSE07.pdf}{Introduction to General Purpose GPU (GPGPU) Programming. ACM Southeast Regional Conference, March 2007.}

\end{itemize}

\section*{Technical expertise}

\begin{itemize}

  \item Languages: Go, C, C++, Python, Rust, Bash, Java, C\#, R

  \item Platforms: Kubernetes, OpenShift, RHACS, Docker, Linux, IBM Power, IBM Z

  \item Security / systems: eBPF, runtime detection, behavioral baselines, policy-as-code, LLM gateways, MCP, secure sandboxes, support automation, agent workflows

  \item Research / methods: symbolic execution, software verification, formal methods, program synthesis, SMT/SAT solvers, LLVM, KLEE

\end{itemize}

\emph{Last updated: June 2026.}

\end{document}
